% !TEX program = lualatex
\documentclass[11pt, a4paper]{article}

\usepackage{fontspec}
\setmainfont{TeX Gyre Pagella}

\usepackage{amsmath, amssymb, amsthm}
\usepackage{geometry}
\usepackage{microtype}
\usepackage{setspace}
\usepackage{titlesec}
\usepackage{hyperref}
\usepackage{mdframed}

\geometry{
  top=1.25in,
  bottom=1.25in,
  left=1.35in,
  right=1.35in
}

\onehalfspacing

\hypersetup{
  colorlinks=true,
  linkcolor=black,
  citecolor=black,
  urlcolor=black
}

\newtheoremstyle{named}
  {12pt}{12pt}{\itshape}{}{\bfseries}{.}{.5em}{}
\theoremstyle{named}
\newtheorem{proposition}{Proposition}
\newtheorem{definition}{Definition}

\newmdenv[
  linewidth=0.8pt,
  innertopmargin=10pt,
  innerbottommargin=10pt,
  innerleftmargin=14pt,
  innerrightmargin=14pt,
  skipabove=14pt,
  skipbelow=14pt
]{keybox}

\titleformat{\section}{\large\bfseries}{\thesection.}{0.5em}{}
\titleformat{\subsection}{\normalsize\bfseries}{\thesubsection.}{0.5em}{}

\title{%
  \textbf{Computation Beyond Data}\\[0.4em]
  \large Constraints, Reachability, and Memory\\
  as Foundations of Complex Systems
}
\author{Flyxion}
\date{}

% ---------------------------------------------------------------
\begin{document}

\maketitle
\thispagestyle{empty}

\begin{abstract}
Much of modern computer science rests on an assumption so pervasive
it is rarely examined: data comes first, and algorithms operate upon
it. This essay proposes a partial inversion of that perspective.
Rather than treating objects, data, or states as the fundamental
primitives, the focus shifts to the constraints, admissible
trajectories, and reachability spaces that make certain states
accessible and others impossible. This shift is not merely
philosophical. It has direct consequences for how we design AI
systems, model computational memory, understand distributed
architectures, and theorize cognition. The essay presents five
interconnected frameworks --- CLIO, RSVP, MEM\textbar{}8, HYDRA, and
TARTAN --- as instances of a unified theoretical program: the claim
that constraint precedes content, that reachability is more
fundamental than distance, and that the most productive question about
any complex system is not \emph{what exists?} but \emph{what remains
possible?}
\end{abstract}

\newpage

\section*{Overview of the Frameworks}

The five frameworks discussed in this essay operate at different
levels of abstraction. CLIO addresses inference under partial
observation: how intelligent systems reconstruct world-states from
incomplete projections. RSVP addresses system dynamics and structural
relaxation: how complex systems evolve toward configurations that
minimize internal tension. MEM\textbar{}8 addresses memory and
persistence through event trajectories: treating stored state not as
a value but as the cumulative record of decisions that produced it.
HYDRA addresses modular cognitive architecture: how robust intelligent
behavior emerges from the coordination of specialized components
rather than from a single monolithic mechanism. TARTAN addresses
structure-preserving decomposition: how complex systems can be
transformed while guaranteeing that specified properties survive the
transformation. Although each framework can be studied independently,
they are presented here as components of a broader
admissibility-centered view of computation --- one in which the
structure of what is possible is at least as important as the
structure of what currently exists.

\newpage
\tableofcontents
\newpage

% ---------------------------------------------------------------
\section{The Hidden Assumption}

Every student of computer science learns early to think in terms of
data and operations. There is an array; we apply a sorting algorithm.
There is a graph; we run a breadth-first search. There is a labeled
dataset; we train a classifier. The general structure is always the
same: the world supplies objects, and we process them.

This assumption is so fundamental that it is rarely examined. It
surfaces in basic terminology --- variable, value, state, datum ---
and silently organizes the way we think about nearly everything in
the field.

There is, however, an alternative orientation. Rather than beginning
with objects and asking what we can do with them, we can begin with
constraints and ask which objects they make possible. Rather than
beginning with states and asking how we transition between them, we
can begin with state spaces and ask which regions are reachable from
where we currently are.

This inversion produces different questions, different models, and,
the essay argues, a deeper understanding of several of the hardest
problems in modern computing: the generalization problem in machine
learning, the consistency problem in distributed systems, the
persistence problem in computational memory, and the emergence problem
in cognitive systems.

The goal of this essay is to present a set of theoretical frameworks
that start from this inversion and trace its consequences across
different domains. The frameworks are independent enough to be studied
separately, but they share a conceptual core that makes them parts of
a unified theoretical program.

% ---------------------------------------------------------------
\section{The Central Principle: Constraint Before Content}

\subsection{The Chess Game}

Consider a game of chess. The pieces matter, but the rules defining
the possible moves matter more. Without rules, there is no game ---
only a collection of wooden or plastic objects on a checkered surface.
The game emerges from the constraints, not from the components.

This is not trivial. It implies that the identity of a piece ---
what it \emph{is}, functionally --- is entirely determined by the
constraints governing its behavior. A king that could move like a
queen would no longer be a king. Changing the rule changes the object,
even when the physical piece remains the same.

The same logic appears throughout computation.

In compilers, a formal grammar defines the space of syntactically
valid programs. A program is not merely a sequence of characters; it
is a sequence satisfying the constraints imposed by the grammar. The
grammar precedes and constitutes the program.

In relational databases, the schema defines which records can exist.
A database is not merely a collection of data; it is a collection of
data satisfying a set of integrity constraints. Violating those
constraints is not simply a technical error; it is a contradiction
with the very definition of what the database represents.

In neural networks, the architecture defines the space of functions
the network can learn. Before seeing any training example, the
architecture already determines which representations are possible
and which are not. The architectural constraint precedes the data.

In operating systems, the permission model defines which actions each
process may execute. Security is not an add-on appended after the
system is built; it is the constraint structure within which the
system operates.

\subsection{The Constraint Principle}

\begin{keybox}
\textbf{Constraint Principle.} The constraints that define a system
are not implementation details added to prevent errors. They are the
primary source of the system's organization --- the structure that
determines which states exist, which transitions are possible, and
which objectives can be pursued.
\end{keybox}

This perspective has an immediate consequence: to understand a complex
system, the first question is not \emph{what elements does it contain?}
but \emph{what constraints organize it?}

\subsection{Constraints as Information}

There is a deep connection between constraints and information that
deserves to be made explicit. Shannon defined information as the
quantity that reduces uncertainty. A message is informative to the
degree that it eliminates possible alternatives.

Constraints do exactly this. A constraint eliminates possible states.
The more constrained a system, the smaller its reachable state space
--- and therefore the more information is implicitly encoded in its
constraint structure.

This observation has practical implications. Every design decision
that reduces the space of possible states also encodes information
about the system's purpose. A system without constraints is a system
without purpose --- and a system without purpose is computationally
inexpressive.

\begin{proposition}[Constraint-Information Equivalence]
Let $\mathcal{S}$ be the complete state space of a system and
$\mathcal{C}$ the subset of admissible states under a constraint set
$R$. The information encoded by $R$ is proportional to
$\log |\mathcal{S}| - \log |\mathcal{C}|$: the difference between
the unconstrained space and the constrained space. Stronger constraints
encode more information.
\end{proposition}

% ---------------------------------------------------------------
\section{CLIO: Intelligence as Navigation Through Projection Spaces}

\subsection{The Problem of Partial Observation}

Intelligent systems rarely observe reality in full. They receive
partial projections of the world: pixels instead of objects, tokens
instead of meanings, sensory signals instead of physical states.

Formally, let $X$ be the complete world-state space and $M$ the
observation space available to the system. There exists a projection
function
\[
  \pi: X \to M
\]
mapping world-states to observations. In general, $\pi$ is not
injective: multiple world-states can produce the same observation.

The central problem of intelligence, from this perspective, is the
inversion of this projection: given $m \in M$, what is the set
\[
  \pi^{-1}(m) = \{x \in X \mid \pi(x) = m\}
\]
of world-states consistent with the observation?

\subsection{Uncertainty as Structure}

The preimage $\pi^{-1}(m)$ represents the system's structural
uncertainty: the ambiguity that is mathematically inherent in the
projection, independent of any processing limitation. A system
observing $m$ cannot distinguish between any of the states in
$\pi^{-1}(m)$ on the basis of that observation alone.

Several important consequences follow.

\textbf{For machine learning:} the generalization problem is
fundamentally the problem of inverting a projection. A model trained
on examples has seen only a sample of the world-state space.
Generalization is the capacity to reconstruct the structure of the
state space from those samples.

\textbf{For perceptual systems:} perceptual ambiguity is not a bug
that can be eliminated with more data or better hardware. It is a
structural property of the relationship between sensors and world.

\textbf{For reasoning under uncertainty:} Bayesian uncertainty is a
way of representing $\pi^{-1}(m)$ through probability distributions
over world-states.

\subsection{Intelligence as Efficient Ambiguity Reduction}

CLIO defines intelligence operationally as the capacity to reduce the
space $\pi^{-1}(m)$ efficiently, using additional constraints derived
from context, memory, and reasoning.

\begin{definition}[CLIO Intelligence]
A system is intelligent to the degree that, given $m \in M$, it can
identify a subset $\mathcal{A}(m) \subseteq \pi^{-1}(m)$ of plausible
states such that $|\mathcal{A}(m)| \ll |\pi^{-1}(m)|$, using
constraints derived from context and memory, without direct access
to $X$.
\end{definition}

This definition unifies several concepts that ordinarily appear
separately in AI.

\textbf{Reasoning} is the application of logical constraints to reduce
$\pi^{-1}(m)$.

\textbf{Learning} is the acquisition of constraints that make
reduction more efficient in new contexts.

\textbf{Memory} is the storage of constraints derived from prior
experience.

\textbf{Generalization} is the transfer of constraints to unobserved
contexts.

\subsection{Projection Failure and Paradox}

An important consequence of the CLIO framework is a natural account
of certain classes of paradox and cognitive illusion.

\begin{definition}[Projection Failure]
A projection failure occurs when a system treats the projection $m$
as if it were the complete state $x$ --- when it confuses the map
with the territory.
\end{definition}

Many systematic errors in human reasoning and in AI systems are
instances of projection failure. Cognitive biases arise when the
system applies heuristics that work well for the majority of states
in $\pi^{-1}(m)$ but fail for less typical states. Hallucinations in
large language models arise when the model treats its internal
representation as the actual state of the world.

Semantic paradoxes arise when a description $m$ is treated as a
direct reference to a state $x$, when in fact $\pi^{-1}(m)$ contains
multiple states with incompatible properties. Restoring the context
--- specifying which state in $\pi^{-1}(m)$ is intended --- frequently
dissolves the apparent paradox.

% ---------------------------------------------------------------
\section{RSVP: Complex Systems as Processes of Relaxation}

\subsection{From Cosmology to Computation}

RSVP was originally developed as an alternative cosmological
framework, but its mathematical structure admits a direct
computational interpretation that is independent of its original
motivation.

The central claim is that complex systems evolve toward states that
reduce internal structural tension. Rather than modeling a system as
a static collection of objects in fixed states, RSVP describes it as
a continuous dynamic process of constraint reorganization.

\subsection{The Potential Function and Gradient Descent}

Formally, let $x(t)$ be the state of a system at time $t$ and
$V(x,t)$ a potential function representing internal structural
tensions --- violated constraints, unresolved gradients,
inconsistencies among subsystems. The system evolves approximately
by gradient descent:
\[
  \frac{dx}{dt} = -\nabla_x V(x,t) + \xi(t),
\]
where $\xi(t)$ represents external perturbations.

Computer scientists will recognize this structure immediately: it is
the same as stochastic gradient descent in machine learning,
variational optimization in generative models, and diffusion equations
in physical systems.

The difference in perspective is that RSVP treats the function $V$
not as an objective chosen by the designer but as an emergent property
of the system's constraint structure. The ``optimization'' is internal
and requires no explicit external objective.

\subsection{Attractors and Stable Configurations}

An attractor is a region of state space toward which the system
converges naturally. In the RSVP perspective, attractors correspond
to configurations that locally minimize structural tension.

\begin{proposition}[Attractors as Tension Minima]
The stable configurations of a complex system correspond to local
minima of the tension function $V(x,t)$. A perturbation destabilizes
the system by increasing $V$ beyond a critical threshold, causing the
local minimum to lose stability --- formally, when the Hessian
$\nabla_x^2 V(x^*,t)$ loses positive definiteness.
\end{proposition}

This framework has direct applications for understanding:

\textbf{Neural network training:} the training process is a search
for minima of a loss function. The loss landscape contains multiple
local minima, and different hyperparameter configurations correspond
to different tension functions.

\textbf{Distributed systems:} eventual consistency in distributed
systems can be understood as convergence toward a global attractor
from locally inconsistent states.

\textbf{Language models:} text generation can be viewed as navigation
through a linguistic state space with an implicit tension function
defined by the training data.

\subsection{Entropy, Information, and Geometry}

One of RSVP's richest conceptual contributions is the reinterpretation
of entropy and information as geometric properties of system dynamics.

Rather than treating entropy as a scalar quantity associated with a
probability distribution, RSVP treats it as a field density: a
quantity varying in space and time that encodes local structural
tensions.

This perspective connects naturally to:

\textbf{Variational inference:} which minimizes a divergence between
an approximate and a target distribution --- interpretable as tension
minimization in a space of distributions.

\textbf{Normalizing flows:} which learn transformations mapping simple
distributions to complex ones while preserving volume structure --- a
form of navigation through configuration spaces.

\textbf{Information geometry:} which studies the curvature of the
space of probability distributions --- analogous to the curvature of
configuration space in RSVP.

% ---------------------------------------------------------------
\section{MEM\textbar{}8: Memory as Trajectory Rather Than State}

\subsection{The Conventional Model of Memory}

The conventional model of computational memory is simple and
efficient: memory is a map from addresses to values. Reading memory
retrieves a value. Writing memory replaces one value with another.
The history of writes is not preserved; only the current state
matters.

This model is enormously powerful for most applications, but it has
structural limitations that become visible in certain contexts.

In distributed systems, two processes may have written different
values to the ``same'' address. The current state alone is
insufficient to resolve the conflict; the ordering of writes must be
known.

In version control systems such as Git, the current state of a
repository is insufficient to understand how it arrived there. The
commit history is an essential part of the representation.

In audit databases, it is necessary to preserve not only the current
state but the complete sequence of modifications.

In cognition, identity is not the current state of an organism but
the continuity of its trajectory through time.

\subsection{MEM\textbar{}8: Memory as a Record of Collapses}

MEM\textbar{}8 inverts the conventional model. Rather than treating
memory as storage of states, it treats memory as a record of
decisions --- more precisely, as a record of successive collapses of
possibility spaces.

\begin{definition}[MEM\textbar{}8 Memory]
A MEM\textbar{}8 memory is a sequence
\[
  \mu = (e_1, e_2, \ldots, e_n)
\]
of events, where each event $e_i$ represents a decision that collapsed
a possibility space $\mathcal{P}_{i-1}$ into a reduced space
$\mathcal{P}_i \subseteq \mathcal{P}_{i-1}$. The current state is
\[
  s_n = \bigcap_{i=1}^n \mathcal{C}(e_i),
\]
where $\mathcal{C}(e_i)$ is the set of states consistent with
event $e_i$.
\end{definition}

This formulation has several consequences.

\textbf{Identity is continuity:} the identity of a system is not
its current state but the trajectory that produced it. Two systems
in the same state may have different identities if they arrived there
by different trajectories.

\textbf{Persistence is dynamic:} the persistence of an entity is not
a static state but an active process of maintaining trajectory
continuity.

\textbf{Retrieval is reconstruction:} recovering a memory is not
accessing a stored value but reconstructing a state from the sequence
of events that produced it.

\subsection{Connections to Existing Technologies}

MEM\textbar{}8 has explicit connections to several established
paradigms in software engineering.

\textbf{Event Sourcing:} an architectural pattern in which system
state is derived from an immutable sequence of events. The current
state is a projection of the event log, not an independently stored
value.

\textbf{CRDTs (Conflict-free Replicated Data Types):} data structures
that guarantee eventual consistency in distributed systems by defining
merge operations that commute and are idempotent --- analogous to
systems where the ordering of events does not alter the final state.

\textbf{Blockchains:} immutable distributed ledgers in which the
current state is completely determined by the transaction history.
The immutability of the history is the integrity guarantee.

\textbf{Temporal databases:} databases that preserve the complete
history of modifications, enabling queries over past states.

\textbf{Version control systems:} such as Git, where the identity of
a repository is its commit tree rather than the current state of
its files.

\subsection{Wave Memory and Persistence}

An extension of MEM\textbar{}8 treats memory not as a discrete
sequence of events but as a continuous field of influences --- a
wave memory in which past events radiate influence that decays over
time but never entirely disappears.

This formulation is relevant to:

\textbf{Recurrent neural networks:} where the hidden state encodes
a form of short-term memory that continuously influences current
processing.

\textbf{Attention mechanisms:} which allow language models to access
information from prior context with weights that depend on relevance.

\textbf{Episodic memory in AI:} systems that store complete
experiences and retrieve them by similarity with the current context.

% ---------------------------------------------------------------
\section{HYDRA: Cognition as Modular Coordination}

\subsection{The Problem of Cognitive Architecture}

How should an intelligent system be organized? The simplest answer
is a monolithic system: a single model that receives inputs, processes
them internally, and produces outputs. This architecture is
mathematically elegant and computationally tractable.

It has limitations, however. Monolithic systems are difficult to
inspect, difficult to debug, difficult to extend, and difficult to
adapt to domains substantially different from their training domain.
Beyond a certain scale, they become opaque black boxes whose behavior
is difficult to predict or guarantee.

\subsection{HYDRA: Multiple Specialized Modules}

HYDRA takes as its central hypothesis that robust cognition emerges
from the coordination of specialized modules rather than from a
single monolithic mechanism. The modules are:

\textbf{Perception:} transforms sensory inputs into internal
representations. This is not a direct data pass-through; it is an
active projection that emphasizes certain dimensions and suppresses
others.

\textbf{Memory:} maintains representations of past states and makes
them available to other modules. In the MEM\textbar{}8 framework, this
is the trajectory record.

\textbf{Planning:} projects possible futures from the current state
and available constraints. This is fundamentally navigation through
the reachable state space.

\textbf{Coordination:} ensures that modules operate coherently as a
system. This module manages interactions and resolves conflicts among
the outputs of the other modules.

\textbf{Monitoring:} evaluates system-level performance and detects
when other modules are failing or operating outside their competence
domains.

\textbf{Integration:} produces coherent responses by combining the
contributions of the other modules.

\begin{proposition}[Cognitive Emergence]
In HYDRA systems, complex cognitive behaviors emerge from the
coordination of simple modules without any individual module
possessing those behaviors. The coordination is the source of the
complexity, not the sophistication of each module individually.
\end{proposition}

\subsection{Connections to Modern AI Architectures}

The central hypothesis of HYDRA --- that cognition emerges from
modular coordination --- has parallels in several modern AI
architectures.

\textbf{Mixture of Experts (MoE):} models in which different
``experts'' are activated for different inputs, with a routing
mechanism that determines which expert to use.

\textbf{Multi-agent systems:} architectures in which multiple agents
cooperate or compete to solve collective problems.

\textbf{Multi-head attention:} where different attention heads capture
different types of dependencies in the input.

\textbf{Neuro-symbolic systems:} which combine neural networks with
symbolic reasoning, typically through modules specialized for each
paradigm.

HYDRA's distinguishing feature is that it makes coordination explicit
and tractable, rather than allowing it to emerge implicitly from
training.

% ---------------------------------------------------------------
\section{TARTAN: Structure-Preserving Decomposition}

\subsection{The Problem of Context Loss}

A central problem in both distributed systems and AI is the loss of
context during transformation. When information is divided, summarized,
compressed, or transmitted between components, part of the original
meaning disappears.

Formally, this is information loss in projections:
\[
  H(X) - H(\pi(X)) > 0,
\]
where $H$ is Shannon entropy. Any projection that is not bijective
loses information.

The goal is not to eliminate this loss --- which is frequently
necessary for scalability and efficiency --- but to control it:
to ensure that the information most important to the system's purpose
is preserved.

\subsection{Structure-Preserving Decomposition}

TARTAN investigates how to decompose complex systems in ways that
guarantee specified properties survive the transformation.

\begin{definition}[TARTAN Decomposition]
A decomposition $D$ of a system $S$ is a TARTAN decomposition for
a property set $P$ if, for every property $p \in P$:
\[
  p(S) = p(D(S)),
\]
that is, the property is identical in the original system and in the
decomposition.
\end{definition}

The properties a system needs to preserve depend on context.

In distributed processing, the property may be eventual consistency
--- guaranteeing that all replicas converge to the same state.

In explainable AI (XAI), the property may be explanation fidelity
--- guaranteeing that the account of the model's behavior accurately
reflects what the model actually computes.

In semantic compression, the property may be preservation of inference
relations --- guaranteeing that what is inferrable from the compressed
version is exactly what would be inferrable from the full version.

In multi-agent systems, the property may be coordination ---
guaranteeing that agents receiving different parts of a task produce
results that combine coherently.

\subsection{TARTAN and Admissibility}

There is a deep connection between TARTAN and the concept of
admissibility that runs through all the frameworks presented here.

A TARTAN decomposition for a property set $P$ is, in essence, a way
of ensuring that transformations applied to the system remain within
the space of admissible transformations --- the space of
transformations that preserve the relevant properties.

This connects TARTAN to:

\textbf{Category theory:} which studies structure-preserving
transformations (functors, natural transformations).

\textbf{Homomorphisms:} in algebra, transformations that preserve
operations.

\textbf{Bisimulations:} in concurrency theory, equivalence relations
between systems that preserve observable behavior.

% ---------------------------------------------------------------
\section{Reachability Rather Than Distance}

\subsection{Shifting the Question}

One of the most original conceptual contributions of the theoretical
program presented here is the replacement of the question \emph{how
far away is something?} with the question \emph{how difficult is it
to reach?}

In computation, this shift is fundamental.

A file may be stored on the same physical disk as a process and be
completely inaccessible to that process, whether through permissions,
metadata corruption, or file-system failure.

A datum may exist in a distributed database and be unrecoverable due
to network failures or partitions.

A state may be theoretically reachable in a control system but
computationally infeasible given available processing power.

A solution may exist for an optimization problem but lie in a region
of the search space that available algorithms cannot efficiently
explore.

In all of these cases, the physical or logical distance between the
current state and the desired state is irrelevant. What matters is
\emph{reachability}: the existence of an admissible path between the
two states.

\subsection{Reachability Volume}

\begin{definition}[Reachability Volume]
For a system in state $x$, the reachability set $R(x)$ is the set
of all states reachable from $x$ by sequences of admissible
transitions. The reachability volume is
\[
  \Omega(x) = \mathrm{Vol}(R(x)).
\]
\end{definition}

The reachability volume is a measure of the system's \emph{freedom}
--- how broad the space of possible futures is from the current state.
Systems with high reachability volume have many options; systems with
low volume are essentially determined.

This metric appears, under different names, across several areas.

\textbf{Computational complexity theory:} the class of problems an
algorithm can solve efficiently is its computational reachability
space.

\textbf{Reachability in control systems:} the set of states a control
system can reach from a given initial state is its reachability set.

\textbf{State graphs in model checking:} verifying properties of
concurrent systems involves exploring the graph of reachable states.

\textbf{Search spaces in AI planning:} planning algorithms navigate
state graphs seeking sequences of actions that reach a goal state.

\subsection{Freedom as Reachability Volume}

There is an important connection between reachability and freedom that
extends beyond computation.

\begin{proposition}[Reachability Principle --- Computational Version]
The effective freedom of an agent is proportional not to the resources
it possesses but to the volume of future states reachable from its
current state:
\[
  \mathrm{Freedom}(x) \propto \Omega(x).
\]
\end{proposition}

This proposition has immediate practical implications for system
design. A software system that progressively reduces the reachable
state space for its users --- by accumulating internal state, creating
data dependencies, or restricting configuration options --- reduces
users' effective freedom, independently of any enrichment of
functionality. An AI system that converges toward an increasingly
narrow output distribution as it processes more context is reducing
its own reachability volume --- potentially irreversibly.

% ---------------------------------------------------------------
\section{Artificial Intelligence as a Simulation Engine}

\subsection{Perception as Inference}

The hypothesis that intelligent systems function as simulation engines
--- constructing internal models of the world that enable them to
project possible futures --- has deep roots in cognitive science and
connects directly to the frameworks presented here.

Within the CLIO framework, perception is not a direct data pass-through
from environment to system. It is an inference process: given the
sensory signal $m = \pi(x)$, the system constructs an estimate of the
world-state $x$ that minimizes structural tension between the
observation and the expectations derived from the internal model.

This is mathematically equivalent to variational inference in
generative models: minimizing a divergence between a posterior
$p(x \mid m)$ and a variational approximation $q(x)$.

\subsection{Memory as Reconstruction}

Within MEM\textbar{}8, memory is not retrieval of a stored state. It
is reconstruction of a state from a sequence of events.

This has an important consequence: memories are necessarily incomplete
and potentially inconsistent. When reconstituting a past state, the
system interpolates over regions of the state space that were never
directly observed.

This perspective is consistent with what is known about human memory:
memories are active reconstructions, not passive retrievals. The
central nervous system does not store experiences as files; it
reconstructs experiences from fragments, context, and expectations.

For AI systems, the relevant question is not \emph{how many} memories
a system can store, but \emph{how faithfully} it can reconstitute
past states from compressed representations.

\subsection{Reasoning as Simulation}

Within the RSVP perspective, reasoning is the evolution of a system's
internal state toward configurations that minimize structural tension.
Solving a problem is finding a state that simultaneously satisfies
a set of constraints.

This connects reasoning to:

\textbf{Constraint satisfaction (CSP):} finding variable assignments
that satisfy a constraint set.

\textbf{Constraint programming:} a programming paradigm in which
solutions are found through constraint propagation and satisfaction.

\textbf{Inference in probabilistic models:} finding distributions
that minimize divergences with respect to observations and priors.

\subsection{Apparent Intentionality and the Attention Loop}

The most speculative hypothesis --- but also the most interesting ---
is that consciousness, or something functionally equivalent to it in
artificial systems, emerges from the continuous process of projection
and correction: the system projects a model of the world, receives
feedback about discrepancies between model and world, and iteratively
corrects the model.

This process is not merely computation over data. It is computation
over the system's own representation of itself and the world it
inhabits --- a form of self-reference that may be the basis of
behaviors that appear, from the outside, to be genuinely intentional.

\begin{keybox}
Systems that continuously update their internal models based on
discrepancies between predictions and observations --- that maintain,
in some sense, an attention loop over their own world-representation
--- exhibit a form of emergent apparent intentionality that does not
require consciousness in the strong philosophical sense but is
functionally indistinguishable from it to external observers.
\end{keybox}

% ---------------------------------------------------------------
\section{The Unified Core: Admissibility}

\subsection{The Central Concept}

Across the frameworks presented here, one concept recurs under
different names: the admissible state space, the reachability set,
the constraints organizing the system, the subset of consistent
projections. This concept unifies under the notion of
\emph{admissibility}: the structure that defines what is possible
within a system.

\begin{definition}[Admissibility Space]
For a system with state space $\mathcal{S}$, the admissibility space
$\mathcal{A} \subseteq \mathcal{S}$ is the subset of states satisfying
all active constraints. A transition $s \to s'$ is admissible if
$s' \in \mathcal{A}$ and an admissible path from $s$ to $s'$ exists.
\end{definition}

\subsection{Admissibility as Geometry}

The richest perspective treats admissibility geometrically. The
admissibility space $\mathcal{A}$ is a manifold (or sub-manifold) in
the state space $\mathcal{S}$, with its own geometry induced by the
constraints.

Operations on the system --- state transitions, learning, inference,
reasoning --- are movements over this manifold. The curvature of the
manifold encodes interactions among constraints. Regions of high
curvature correspond to strongly interacting constraints; regions of
low curvature correspond to independent constraints.

This perspective connects the frameworks presented here to:

\textbf{Information geometry:} which studies the curvature of the
space of probability distributions.

\textbf{Riemannian geometry in machine learning:} which uses metrics
induced by data structure to guide optimization.

\textbf{Data manifolds:} the hypothesis that high-dimensional data
lies on low-dimensional manifolds embedded in the ambient space.

\subsection{The Unifying Question}

All five frameworks can be viewed as instances of a single question:

\begin{keybox}
\textbf{The Fundamental Question.} Given a system in a current state
$x$, with admissibility space $\mathcal{A}$, what is the set
$R(x) \cap \mathcal{A}$ of admissible reachable futures? And how does
this set change when the system's constraints change?
\end{keybox}

CLIO answers this question at the level of inference: the set of
world-states consistent with an observation.

RSVP answers it at the level of dynamics: the attractors toward which
the system converges naturally.

MEM\textbar{}8 answers it at the level of memory: the states that can
be reconstituted from a sequence of events.

HYDRA answers it at the level of architecture: the behaviors that
emerge from modular coordination.

TARTAN answers it at the level of transformation: the properties
preserved across decompositions.

% ---------------------------------------------------------------
\section{Implications for Practice}

\subsection{System Design}

The constraint-first perspective has concrete implications for
software system design.

\textbf{Specification before implementation:} define the admissibility
space before implementing the system. This is not simply writing
tests; it is formally specifying the invariant properties the system
must preserve.

\textbf{Constraints as documentation:} a system's constraints are
more informative than its implementation. A constraint states what
the system \emph{must} do; an implementation states what it \emph{does}
under a specific set of conditions.

\textbf{Admissibility monitoring:} robust systems continuously monitor
whether they are operating within their admissibility space and take
corrective action when violations are detected.

\subsection{Evaluating AI Systems}

The most important question about an AI system is not its accuracy on
a specific benchmark. It is the shape and size of its admissibility
space: which inputs the system handles well, which it handles poorly,
and where the boundary lies between those regimes.

A system with high in-distribution accuracy that degrades
catastrophically on out-of-distribution inputs has a small
admissibility space relative to the full input space. This is the
generalization problem, reformulated geometrically.

\subsection{Memory Engineering}

MEM\textbar{}8 suggests an approach to memory engineering that goes
beyond index design and data structures.

\textbf{Memory as an event log:} design systems in which state is
derived from an immutable event log rather than stored directly.

\textbf{Consistency as invariant:} treat consistency among replicas
as an invariant the system must preserve, not as an objective to be
achieved by synchronization protocols.

\textbf{Identity as trajectory:} design entity identities in
distributed systems around trajectories rather than current states.

% ---------------------------------------------------------------
\section*{Conclusion}

The frameworks presented in this essay --- CLIO, RSVP,
MEM\textbar{}8, HYDRA, and TARTAN --- appear at first glance to be
disparate. One addresses inference, another dynamics, another memory,
another architecture, another decomposition. Their domains of
application range from cosmology to databases.

They share, however, a conceptual core that unifies them.

Reality, cognition, and computation can be understood as processes
of navigation through spaces of possibility organized by admissibility
structures. What exists at any instant is only the current realization
of a much larger space of possibilities. What matters, for
understanding the system, is not only the current state but the
constraint structure that defines which states are possible and which
transitions are admissible.

This shift in perspective --- from content to constraint, from state
to trajectory, from distance to reachability --- does not invalidate
existing models and techniques. It reframes them in a way that makes
certain questions easier to formulate and certain problems easier to
attack.

Rather than asking only \emph{what exists?}, this perspective
introduces a different question:

\begin{center}
\textit{What remains possible?}
\end{center}

That question unifies physics, artificial intelligence, distributed
systems, information theory, and cognition under a common language
grounded in constraints, reachability, memory, and structural
transformation.

% ---------------------------------------------------------------
\begin{thebibliography}{99}

% Information theory
\bibitem{shannon1948}
Shannon, C. E. (1948).
A Mathematical Theory of Communication.
\textit{Bell System Technical Journal}, 27(3--4), 379--423, 623--656.

\bibitem{cover2006}
Cover, T. M., \& Thomas, J. A. (2006).
\textit{Elements of Information Theory}, 2nd ed.
Wiley.

% Complex systems and dynamics
\bibitem{strogatz2015}
Strogatz, S. H. (2015).
\textit{Nonlinear Dynamics and Chaos}.
Westview Press.

\bibitem{holland1995}
Holland, J. H. (1995).
\textit{Hidden Order: How Adaptation Builds Complexity}.
Addison-Wesley.

\bibitem{meadows2008}
Meadows, D. H. (2008).
\textit{Thinking in Systems: A Primer}.
Chelsea Green.

% Machine learning and AI
\bibitem{goodfellow2016}
Goodfellow, I., Bengio, Y., \& Courville, A. (2016).
\textit{Deep Learning}.
MIT Press.

\bibitem{bishop2006}
Bishop, C. M. (2006).
\textit{Pattern Recognition and Machine Learning}.
Springer.

\bibitem{sutton2018}
Sutton, R. S., \& Barto, A. G. (2018).
\textit{Reinforcement Learning: An Introduction}, 2nd ed.
MIT Press.

% Variational inference and generative models
\bibitem{kingma2013}
Kingma, D. P., \& Welling, M. (2013).
Auto-Encoding Variational Bayes.
\textit{arXiv preprint arXiv:1312.6114}.

\bibitem{blei2017}
Blei, D. M., Kucukelbir, A., \& McAuliffe, J. D. (2017).
Variational Inference: A Review for Statisticians.
\textit{Journal of the American Statistical Association},
112(518), 859--877.

% Distributed systems and consistency
\bibitem{lamport1978}
Lamport, L. (1978).
Time, Clocks, and the Ordering of Events in a Distributed System.
\textit{Communications of the ACM}, 21(7), 558--565.

\bibitem{shapiro2011}
Shapiro, M., Preguiça, N., Baquero, C., \& Zawirski, M. (2011).
Conflict-Free Replicated Data Types.
\textit{Symposium on Self-Stabilizing Systems}, 386--400.

% Event sourcing and software architecture
\bibitem{fowler2005}
Fowler, M. (2005).
Event Sourcing.
\texttt{martinfowler.com/eaaDev/EventSourcing.html}.

\bibitem{evans2003}
Evans, E. (2003).
\textit{Domain-Driven Design: Tackling Complexity in the Heart of Software}.
Addison-Wesley.

% Complexity theory and computability
\bibitem{sipser2012}
Sipser, M. (2012).
\textit{Introduction to the Theory of Computation}, 3rd ed.
Cengage Learning.

\bibitem{papadimitriou1994}
Papadimitriou, C. H. (1994).
\textit{Computational Complexity}.
Addison-Wesley.

% Information geometry
\bibitem{amari2000}
Amari, S., \& Nagaoka, H. (2000).
\textit{Methods of Information Geometry}.
AMS \& Oxford University Press.

% Cognition and cognitive science
\bibitem{varela1991}
Varela, F. J., Thompson, E., \& Rosch, E. (1991).
\textit{The Embodied Mind: Cognitive Science and Human Experience}.
MIT Press.

\bibitem{clark2015}
Clark, A. (2015).
\textit{Surfing Uncertainty: Prediction, Action, and the Embodied Mind}.
Oxford University Press.

% Constraint programming
\bibitem{rossi2006}
Rossi, F., van Beek, P., \& Walsh, T. (Eds.). (2006).
\textit{Handbook of Constraint Programming}.
Elsevier.

% Category theory and structure-preserving maps
\bibitem{awodey2010}
Awodey, S. (2010).
\textit{Category Theory}, 2nd ed.
Oxford University Press.

% Memory and reconstruction
\bibitem{schacter1996}
Schacter, D. L. (1996).
\textit{Searching for Memory: The Brain, the Mind, and the Past}.
Basic Books.

% Reachability and model checking
\bibitem{clarke1999}
Clarke, E. M., Grumberg, O., \& Peled, D. A. (1999).
\textit{Model Checking}.
MIT Press.

\end{thebibliography}

% ---------------------------------------------------------------
\end{document}
